\documentclass[a4paper,12pt]{report}
\renewcommand\thesection{\arabic{section}}
\renewcommand{\contentsname}{Cuprins}
\renewcommand{\figurename}{Figura}
\renewcommand{\tablename}{Tabel}
\usepackage{amsmath}
\usepackage{rotating}
\usepackage{listings}
\usepackage{lipsum}  


\setcounter{tocdepth}{6}
\setcounter{secnumdepth}{6}

\usepackage{url}
\usepackage[backend=bibtex,style=numeric,citestyle=numeric,backref=true,sorting=none]{biblatex}
\addbibresource{bibliography.bib}

\usepackage{indentfirst}
\usepackage[nobottomtitles*]{titlesec}

\newcommand*{\justifyheading}{\raggedleft}
\titleformat{\section}{\normalfont\Large\justifyheading}{\thesection}{18pt}{}

\usepackage{graphicx}
\usepackage[export]{adjustbox}
\graphicspath{ {./images/} }

\usepackage[center]{caption}
\usepackage{subcaption}

\usepackage[margin=3.9cm]{geometry}


\usepackage{listings}
\usepackage{color}
\usepackage{float}

\definecolor{codegreen}{rgb}{0,0.6,0}
\definecolor{codegray}{rgb}{0.5,0.5,0.5}
\definecolor{codepurple}{rgb}{0.58,0,0.82}
\definecolor{backcolour}{rgb}{0.96,0.96,0.96}

\lstdefinestyle{mystyle}{
	backgroundcolor=\color{backcolour},   
	commentstyle=\color{codegreen},
	keywordstyle=\color{blue},
	numberstyle=\tiny\color{codegray},
	stringstyle=\color{codepurple},
	basicstyle=\fontsize{10}{10}\selectfont\ttfamily,
	breakatwhitespace=false,         
	breaklines=true,                 
	captionpos=b,                    
	keepspaces=true,                 
	numbers=left,                    
	numbersep=2pt,                  
	showspaces=false,                
	showstringspaces=false,
	showtabs=false,                  
	tabsize=2
}

\lstset{style=mystyle}

\usepackage{pythonhighlight}

\usepackage{fontspec}% -> Schimbat compilator in LuaLatex
\setmainfont{[UTSans-Regular.otf]}
\usepackage{setspace}
\usepackage{tocloft}
\setstretch{1} 

\begin{document}
	
\include{./cover/cover}
\include{./cover/innercover}
\newpage
\setcounter{page}{2}
\tableofcontents
\newpage
\pagenumbering{arabic}
\renewcommand{\bibname}{Bibliografie}
\renewcommand\lstlistlistingname{\normalsize{CODURI SURSĂ}}
\renewcommand{\lstlistingname}{\normalsize{Secvență de Cod}}
\renewcommand{\listfigurename}{FIGURI}
\renewcommand{\cftloftitlefont}{\normalsize}
\renewcommand{\cftfigfont}{\normalsize}
\renewcommand{\cftfigpagefont}{\normalsize}
\renewcommand{\listtablename}{TABELE}
\renewcommand{\cftlottitlefont}{\normalsize}
\renewcommand{\cfttabfont}{\normalsize}
\renewcommand{\cfttabpagefont}{\normalsize}

% acestea sunt actualizate in mod automat, folosind table, figure si lstlisting
\section*{Lista de figuri, tabele și coduri sursă}
\addcontentsline{toc}{section}{Lista de figuri, tabele și coduri sursă}

\begingroup
\let\clearpage\relax
\listoffigures
\let\clearpage\relax
\listoftables
\let\clearpage\relax
\lstlistoflistings
\setcounter{page}{5}
\endgroup


\clearpage

% trebuie completata manual
\section*{Lista de acronime}
\addcontentsline{toc}{section}{Lista de acronime}

% (se vor scrie toate prescurtările utilizate in ordine alfabetică)

\noindent 
AI - Artificial Intelligence; \\
ASCII - American Standard Code for Information Interchange; \\

\clearpage


\section{Sectiune}

\lipsum[2-4]
% folositi clear page pentru a incepe fiecare sectiune pe o pagina noua 
\clearpage


\section{Alta sectiune}

\subsection{Subsectiune}

\lipsum[2-4]

\subsection{Alta subsectiune}

Citarea se ralizeaza folosind cite~\cite{ex}.

\clearpage


\section{Figuri}

\begin{figure}[htbp]
\centerline{\includegraphics[width=\columnwidth,keepaspectratio]{images/top.png}}
\caption[Exemplu de figura. (titlu in lista de figuri)]{Exemplu de figura. (Explicatia pozei)}
\label{fig:top}
\end{figure}

\clearpage


\section{Fragmente de cod}

\begin{lstlisting}[language=Python,caption={Fragment de cod exemplu},captionpos=b,label={code:sobel}]
print("Hello world!")
\end{lstlisting}

\clearpage


\section{Tabel}

\begin{table}[htbp]
    \centering
    \resizebox{\columnwidth}{!}{
    \begin{tabular}{|c|c|c|c|}
        \hline
         Etapă & Core i7-4710HQ {[}\%{]} & SOC-FPGA {[}\%{]} & 4{\times}ARM A53@1.5MHz {[}\%{]} \\ [0.5ex]\hline\hline
         Detecție           & 91    & 14,20 & 94,4 \\ \hline
         Urmărire           & 0,5   & 0,54  & 2,5 \\ \hline
         Așteptare cadru    & 6     & 81,15 & 2,5 \\ \hline
         Altele             & 3,5   & 4,1   & 1,23 \\ \hline
    \end{tabular}}
    \caption[Exemplu tabel. (text pentru lista de tabele)]{Exemplu tabel. (explicatii pentru tabel)}
    \label{table:comp}
\end{table}

\clearpage
        
\printbibliography
\addcontentsline{toc}{section}{Bibliografie}
\clearpage


\section*{Rezumat}
\addcontentsline{toc}{section}{Rezumat}

(o descriere sumară, în limba română, a principalelor aspecte tratate în proiect – maxim 300 de cuvinte)

\clearpage


\section*{Abstract}
\addcontentsline{toc}{section}{Abstract}

(short description, in english, of main aspects treated in the project)

\clearpage

\end{document}